\documentclass{article}


\usepackage[left=0.5in, right=0.5in, top=1in, bottom=1in]{geometry} % Specific margins
\usepackage{tabularx}
\usepackage{amsmath}
\usepackage{amssymb}



\begin{document}
 \hfill \textbf{\huge Some proofs}   \hfill  \textbf{\large 
	Chad } 
	\rule{19.5cm}{0.8pt} 


	The kaTex rendering for\verb|.md| files is horrendus, 
	thus here is a nicer \LaTeX \ format 

	
\section{Simple Linear Regresion model}
	
	We define our SLR model. 
\begin{equation}
 Y_i = \beta_0 + \beta_1X_i + \epsilon_i
\end{equation}
	
\noindent Where $Y_i$ is our response variable, $\beta_1$ and $\beta_0$ is both 
our parameters and $\epsilon_i$ is the random error term.
We can start off with our assummptions for SLR, 
which are $\epsilon_i \overset{i.i.d}{\sim} N(0, \sigma^2)$. Since $\epsilon_i$ is a
random variable with an assumed Gaussian distribution, It is obvious to see that $Y_i$ 
is then a random variable from the Gaussian distribution. Such that
\begin{equation} \tag{1.2}
Y \sim N(\beta_0 + \beta_1X_i, \sigma^2)
\end{equation}

\noindent To find this is trival, we take the expected value of $Y_i$ such that 
$$\mathbb{E}[Y_i]= \mathbb{E}[\beta_0 + \beta_1X_i + \epsilon_i]
= \mathbb{E}[\beta_0] + \mathbb{E}[\beta_1X_i] + \mathbb{E}[\epsilon_i] = 
\beta_0 + \beta_1X_i
$$
\begin{equation} \tag{1.3}
	\mathbb{E}[Y_i] = \beta_0 + \beta_1X_i 
\end{equation}
$X_i$ is our known constant, $\beta_0$ and $\beta_1$ are the unknown constants, and
$\epsilon_i$ is assumed to have $\mathbb{E}[\epsilon_i] = 0$. 

\subsection{Least Squares method} \tag{1.4}
	
\noindent We are sometimes interested in the observed values deviation 
	from its mean to distinguish the variabillity 
	within the distribution of $Y_i$. Let $Q$ define the least squares equation 
	for $Y_i$.

\begin{equation} \tag{1.4.1}
		Q = \sum_{i=1}^n (Y_i - \mathbb{E}[Y_i])^2
\end{equation}
A high $Q$ indicates high variability within the observed values, this 
is a cause of problems for when building a model to draw predictions. We 
can minimize $Q$ w.r.t the regression model parameters. In order to do this
we borrow the equation from $(1.3)$ and set the partial derivative equal to
zero.
\begin{equation} \tag{1.5.1}
	\frac{\partial Q}{\partial \beta_1} = -2\sum_{i=1}^n X_i(Y_i - 
	\beta_0-\beta_1X_i) = 0
\end{equation} \\
\begin{equation} \tag{1.5.2}
\frac{\partial Q}{\partial \beta_0} = -2\sum_{i=1}^n (Y_i - 
	\beta_0-\beta_1X_i) = 0 
\end{equation}

\noindent We apply the chain rule. From the above equations, we can
derive the normal equations and the linear estimators $b_1$ and $b_0$. We
will conceal the index of the sums from here on out. Below is the first 
normal equation.
\begin{align*}
	  -2\sum X_i(Y_i - \beta_0-\beta_1X_i) = 0 \\
	 \sum (Y_iX_i - b_0X_i - b_1X_i^2) = 0  \\
	 \sum Y_iX_i - b_0\sum X_i - b_1\sum X_i^2 =  0  \\ \\
	\sum Y_iX_i = b_0\sum X_i + b_1 \sum X_i^2 \quad \quad (1.6.1)  
\end{align*}

\begin{align*}
	  -2\sum (Y_i - \beta_0-\beta_1X_i) = 0 \\
	 \sum (Y_i - b_0 - b_1X_i) = 0  \\
	 \sum Y_i - nb_0 - b_1\sum X_i =  0  \\ \\
	\sum Y_i = nb_0 + b_1 \sum X_i \quad \quad (1.6.2)  
\end{align*}

(1.6.1) & (1.6.2) are the derived normal equations.

\section{Least Square Estimators}

	Directly from our normal equation we can derive ourlinear estimators, 
	which by the Gauss-Markov Theorem, OLS estimators are BLUE - best linear 
	unbiased estimators, among all linear ubiased estimators. 
	
   \subsection{Deriving our linear estimators}
	\begin{align*}
		\sum Y_i &= nb_0 + b_1 \sum X_i \quad \quad (\text{rearrange the
		equation.}) \\
			b_0 &= \frac{1}{n}(\sum Y_i -b_1\sum X_i) \\ 
			b_0 &= \bar{Y} - b_1\bar{X}
	\end{align*}

	We can use this equation for $b_0$ to find $b_1$
	\begin{align*}
		\sum Y_iX_i &= b_0\sum X_i + b_1 \sum X_i^2 \quad \quad (\text{
			replace $b_0$ with its equation.}) \\
		\sum Y_iX_i &= (\frac{1}{n}\sum Y_i - b_1\frac{1}{n}\sum X_i)
		\sum X_i + b_1 \sum X_i^2 \\ 
		\sum Y_iX_i &= \frac{1}{n}\sum Y_iX_i - b_1\frac{1}{n}\sum X_i^2
		 + b_1 \sum X_i^2 \\ 
		\sum Y_iX_i - \frac{1}{n}\sum Y_iX_i &= - b_1\bigg(\frac{1}{n}
		\sum X_i^2 - \sum X_i^2\bigg) \quad \quad (\text{
			rearrange the equation.}) \\
		b_1 &=\frac{\sum Y_iX_i - \frac{1}{n}\sum Y_iX_i}{\frac{1}{n} 
		\sum X_i^2 - \sum X_i^2}
	\end{align*}






		






\end{document}
